\documentclass[tikz,border=6pt]{standalone}
\usepackage{ctex}
\usepackage{amsmath}
\usetikzlibrary{calc, arrows.meta}

\begin{document}
\begin{tikzpicture}[scale=1.1, thick]

% 推导示意图：距离相等的几何意义
% 抛物线 y^2 = 4x，焦点 F(1,0)，准线 x=-1
% 取三个不同位置的点，展示 |PF| = |PA|

% 坐标轴
\draw[->, thin, gray] (-2.0, 0) -- (5.0, 0) node[right, font=\small] {$x$};
\draw[->, thin, gray] (0, -3.2) -- (0, 3.2) node[above, font=\small] {$y$};
\node[font=\footnotesize, below right] at (0,0) {$O$};

% 抛物线
\draw[black, thick, domain=-2.9:2.9, smooth, samples=80]
  plot ({(\x)^2/4}, \x);

% 准线 x = -1
\draw[red, thick] (-1, -3.0) -- (-1, 3.0);
\node[red, font=\small, above left] at (-1, 3.0) {准线 $l: x=-\dfrac{p}{2}$};

% 焦点 F
\fill[black] (1, 0) circle [radius=2.5pt];
\node[below right, font=\small] at (1, 0.05) {$F\!\left(\dfrac{p}{2},0\right)$};

% 顶点
\fill[black] (0, 0) circle [radius=2pt];

% === 点 P1（上方，参数 t=2）===
\coordinate (P1) at (1, 2);
\coordinate (A1) at (-1, 2);
\fill[blue] (P1) circle [radius=2.5pt];
\fill[blue] (A1) circle [radius=2pt];
\node[right, font=\footnotesize, blue] at (P1) {$P_1$};
\node[left, font=\footnotesize, blue] at (A1) {$A_1$};
\draw[blue, dashed, thick] (P1) -- (1,0); % P1F
\draw[blue, thick] (P1) -- (A1); % P1A1
% 直角标记
\draw[blue, thin] (-1+0.15, 2) -- (-1+0.15, 2-0.15) -- (-1, 2-0.15);

% === 点 P2（中间，参数 t=0，顶点）===
% 顶点到准线距离 = 1，到焦点距离 = 1
\coordinate (P2) at (0, 0);
\coordinate (A2) at (-1, 0);
\fill[black] (A2) circle [radius=2pt];
\node[below, font=\footnotesize] at (A2) {$A_2$};
\draw[gray, dashed] (P2) -- (1,0);   % P2F = 1
\draw[gray] (P2) -- (A2);             % P2A2 = 1
% 标注
\node[font=\footnotesize, above] at (0.5, 0.15) {$\frac{p}{2}$};
\node[font=\footnotesize, above] at (-0.5, 0.15) {$\frac{p}{2}$};

% === 点 P3（下方，参数 t=-2.4）===
\coordinate (P3) at ({(-2.4)^2/4}, -2.4); % (1.44, -2.4)
\coordinate (A3) at (-1, -2.4);
\fill[blue] (P3) circle [radius=2.5pt];
\fill[blue] (A3) circle [radius=2pt];
\node[right, font=\footnotesize, blue] at (P3) {$P_3$};
\node[left, font=\footnotesize, blue] at (A3) {$A_3$};
\draw[blue, dashed, thick] (P3) -- (1,0); % P3F
\draw[blue, thick] (P3) -- (A3); % P3A3
\draw[blue, thin] (-1+0.15, -2.4) -- (-1+0.15, -2.4+0.15) -- (-1, -2.4+0.15);

% 底部核心结论
\node[font=\small, align=center] at (1.5, -3.8)
  {任意 $P$ 在抛物线上 $\Leftrightarrow$ $|PF| = |PA|$（$A$ 为 $P$ 在准线上的射影）};

\end{tikzpicture}
\end{document}
